% s13_twk_proof: the proof of thm:s13:TWK (input from sec/s13.tex)
\begin{proof}
Throughout, $\cF=(f_1,\dots,f_n;h_1,\dots,h_p)$ is a family of real sources with supports in $Q_R$, $J\ge J_1(R)$, so that $MJ\ge R+\ell_c\ge R+10$: every source lies at $\ell^\infty$-distance $>1$ from $\R^3\setminus\CJ$, the support condition \eqref{eq:s04:S} holds for $\XJ$, and $J\ge J_0(\cF)$. For $\zeta\in\Dtube$ write $g=g(\zeta)$ and $k=k(\zeta)=k_+-k_-$ (\cref{def:s02:tube}); by \cref{lem:s02:E}(E3) the sampled sources satisfy $0\le2k_+\le s_1$ and $0\le2k_-\le H$ on $\Lk$. The inputs are: \Itwo\ with amplitude $s_1\in(0,\sigma_1)$ and constant $C_U$, from \Tthree\ and the torus inputs \cref{hyp:s03:T1,hyp:s03:T2} through \cref{prop:s10:U}, for $a\ne z$ with the constant $C_U+s_1\sup_k|z_k-a_k|$, finite by \cref{def:s03:adm} (\cref{hyp:s12:SCH}); \Ithree, which is \cref{prop:s10:RP}; \Ifour, which is \cref{lem:s02:BFS}; and the cell constant $B^{(a)}_{\rm cell}<\infty$ of the comparison convention (\cref{hyp:s12:Bcell}). Under these, \cref{cor:s11:W2prime} and \cref{prop:s11:W8} hold on the comparison convention and \cref{lem:s04:E}(d) on King's.

\emph{Step 1 (modulus).} \Cref{lem:s02:alg}(d) holds for every state of \cref{def:s02:acube}, killed sites included: $|Z^{a,K}_{k,J}(\zeta)|\le Z^{a,K}_{k,J}(\Re\zeta)$. It therefore suffices to bound $Z^{a,K}_{k,J}$ at the real point $\Re\zeta$, where $Z^{a,K}_{k,J}(\Re\zeta)=N^{a,K}(g,k)\defeq e^{-a_k\sS(k)}\avK{e^{F(g)+\Wick k}}$.

\emph{Step 2 (curvature off).} Let $\langle\cdot\rangle_{K,k}$ be King's cube with each single-site measure at an alive site multiplied by $e^{\eps^3k(x)r_x^2}$; it is a BFS measure (the factor keeps $\int e^{\alpha r^2}d\rho_x<\infty$ because $\lambda>0$, and the killed sites are unchanged), and $N^{a,K}(g,k)=N^{a,K}(0,k)\,\langle e^{F(g)}\rangle_{K,k}$. Its Gaussian is the Neumann photon Gaussian of $\XJ$, the same for both conventions, with covariance $C$ as in \cref{lem:s05:C1}. By \eqref{eq:s04:S} every plaquette with $g\ne0$ has its corners in $\XJ$, so $F(g)=\ip vA_k$ with $v=\dstar_kg$, and \cref{lem:s02:GD} (with $s=1$) and \cref{lem:s05:C1} give $\langle e^{F(g)}\rangle_{K,k}\le e^{\ip v{Cv}_k/2}\le e^{\norm g_{2,k}^2}$. This is the chain of \cref{thm:s05:CK}(a) in the presence of the real scalar source $k$. Hence $N^{a,K}(g,k)\le e^{\norm g_{2,k}^2}N^{a,K}(0,k)$.

\emph{Step 3 (the split).} \Cref{lem:s02:alg}(f) at $f=0$ with $k_1=k_+$, $k_2=k_-$: $N^{a,K}(0,k_+-k_-)\le N^{a,K}(0,2k_+)^{1/2}N^{a,K}(0,-2k_-)^{1/2}$.

\emph{Step 4 (the positive leg on King's convention).} Put $q\defeq2k_+$, a lattice function with $0\le q\le s_1$ supported in $Q_R\cap\Lk\subset\alive^K_J$. The observable $e^{\sum_x\eps^3q(x)r_x^2}$ is $\Theta_{0,q'}$ with $q'=\eps^3q\ge0$, of class $\cI$, and \eqref{eq:s04:S} holds for $\XJ$; Step~1 of the proof of \cref{lem:s04:G} (the kill step on $\cI$) gives $\avK{\Theta_{0,q'}}\le\ava{\Theta_{0,q'}}$, and the multiplier $e^{-(a_k+2\sW_\eps)\sS(q)}$ is one positive constant for both states (\cref{lem:s02:alg}(a)). Hence
\[
N^{a,K}_{k,\CJ}(0,q)\ \le\ N^a_{k,\CJ}(0,q)\ \le\ e^{\kappa\norm q_{1,k}}
\]
by \cref{cor:s11:W2prime}, whose hypotheses \Itwo, \Ithree\ hold (this is the cone instance of \cref{rem:s04:F}, as in \cref{rem:s11:Kconvention}). So $\frac12\ln N^{a,K}(0,2k_+)\le\kappa\norm{k_+}_{1,k}$.

\emph{Step 5 (the negative leg on King's convention).} Put $q\defeq2k_-$, a lattice function with $0\le q\le H$ supported in $Q_R$, and $J\ge J_1(R)$. \Cref{lem:s04:E}(d) gives $\ln N^{a,K}_{k,\CJ}(0,-q)\le K_-\wnorm q$ with the printed constant $K_-$ of \eqref{eq:s11:constants}, built from the cell constant $B^{(a)}_{\rm cell}$ of the comparison cell, which is finite by \cref{hyp:s12:Bcell}: by \cref{lem:s04:E}(c) every cell on which one of the cutoff sources $8\chi_sq$ does not vanish is a full cell inside $\CJ$ containing no killed site, so each factor of the wall insertion of \cref{lem:s04:E}(a) that carries a source is a comparison-cell functional, to which \cref{lem:s11:W4,lem:s11:W5} apply as printed, and the sum over cells is \cref{lem:s11:W67}; the scheme multiplier factorises over the cells (\cref{lem:s02:alg}(a)). The direction of the kill matters here and is the content of \cref{lem:s04:E}(e). Hence $\frac12\ln N^{a,K}(0,-2k_-)\le\frac12K_-\wnorm{2k_-}=K_-\wnorm{k_-}$.

\emph{Step 6 (assembly).} Steps 1--5 give, for every $k\ge0$, $J\ge J_1(R)$ and $\zeta\in\Dtube$,
\[
|Z^{a,K}_{k,J}(\zeta)|\le Z^{a,K}_{k,J}(\Re\zeta)\le\exp\big(\norm{g(\zeta)}_{2,k}^2+\kappa\norm{k_+(\zeta)}_{1,k}+K_-\wnorm{k_-(\zeta)}\big),
\]
which is \eqref{eq:s13:WK}. The constants are those of \eqref{eq:s11:constants}, with $B^{(a)}_{\rm cell}$ in $K_-$ and, for $a\ne z$, $C_U+s_1\sup_k|z_k-a_k|$ in $\kappa$.

\emph{Step 7 (the fixed-cube limits).} Each $Z^{a,K}_{k,J}$ is entire in $\zeta\in\C^{n+p}$: it is a finite-dimensional integral with Gaussian decay in $A$ ($\lmu>0$) and $e^{-\lambda\eps^3r^4}$ decay in $r$ at every alive site, the killed sites carrying no variables, and differentiation under the integral sign is justified. On a compact $Q\subset\Dtube$ the family $\{Z^{a,K}_{k,J}\}_{k\ge0}$ is uniformly bounded, by \eqref{eq:s13:WK} and \cref{lem:s02:E}(E1)--(E3): $\norm g_{2,k}^2\le C_0\fournorm g^2$, $\norm{k_+}_{1,k}\le C_0'\wnorm{k_+}$ and $\wnorm{k_\pm}\le\sum_j|\Re w_j|\wnorm{h_j}$, all continuous in $\zeta$. At every real point of $\Dtube$ the sequence converges: there $2\norm{k_+}_\infty<\sigma_+=s_1=s_W$ and the negative part is unrestricted, so \cref{hyp:s03:W1}, which is posed on King's convention (\cref{cor:k02:W1}(a)), gives the limit in $(0,\infty)$ for the reference scheme; for the admissible scheme $a$, $Z^{a,K}_{k,J}=e^{-(a_k-z_k)\sum_jw_j\sS_k(h_j)}Z^{z,K}_{k,J}$ with $a_k-z_k$ convergent and $\sS_k(h_j)\to\int h_j$ (\cref{def:s03:adm}, \cref{rem:k02:scheme}). The real points of $\Dtube$ form a non-empty open subset of $\R^{n+p}$ (\cref{def:s02:tube}), a set of uniqueness in the connected open $\Dtube$; by Montel's and Vitali's theorems (the fact (T3) stated before \cref{thm:s11:W}),$Z^{a,K}_{k,J}\to G^K_J$ locally uniformly on $\Dtube$, $G^K_J$ is holomorphic there, and \eqref{eq:s13:WK} passes to the limit with $\norm g_{2,k}\to\norm g_2$, $\norm{k_+}_{1,k}\to\norm{k_+}_{L^1}$ ($k_+$ is continuous with compact support) and the sampled $\wnorm{k_-}$ at most the continuum one (E3):
\begin{equation}\label{eq:s13:WKlim}
|G^K_J(\zeta)|\le\exp\big(\norm{g(\zeta)}_2^2+\kappa\norm{k_+(\zeta)}_{L^1}+K_-\wnorm{k_-(\zeta)}\big)\qquad(\zeta\in\Dtube,\ J\ge J_1(R)),
\end{equation}
a bound free of $J$. At real points $G^K_J>0$ (\cref{hyp:s03:W1}).

\emph{Step 8 (the cube limit on a sub-tube).} Split $h_j=h_j^+-h_j^-$ with $h_j^-\defeq\norm{h_j}_\infty\chi_j$, $\chi_j\in C^\infty_c$, $0\le\chi_j\le1$, $\chi_j=1$ on $\supp h_j$, $\supp\chi_j\subset Q_{R+1/2}$, and $h_j^+\defeq h_j+h_j^-\ge0$. For the split family $\cF''=(f;h^+,h^-)$ put $F^K_{k,J}(z,w^+,w^-)\defeq Z^{a,K}_{k,J}(\sum_iz_if_i,\sum_j(w_j^+h_j^++w_j^-h_j^-))$ and $\cD''\defeq\{(z,w^+,w^-)\in\C^{n+2p}:|\Re w_j^\pm|<\delta_j\}$ with $\sum_j\delta_j(\norm{h_j^+}_\infty+\norm{h_j^-}_\infty)<H/2$. On $\cD''$ the scalar source $k''$ of the split family has $\norm{k''_\pm}_\infty\le\norm{k''}_\infty<H/2\le\sigma_+/2$, so $\cD''$ lies in the tube of $\cF''$, whose supports lie in $Q_{R+1/2}$; Steps 1--7 for $\cF''$ and $J\ge J_1(R+\frac12)$ give $F^K_{k,J}\to F^K_J$ locally uniformly on $\cD''$ with the $J$-free bound \eqref{eq:s13:WKlim} for $\cF''$. On $\Omega''_+\defeq\cD''\cap((\mathrm i\R)^n\times(0,\infty)^{2p})$ the curvature coefficients are imaginary and the total scalar source $\sum_j(w_j^+h_j^++w_j^-h_j^-)$ is nonnegative, so the observable is of class $\cI$ and \eqref{eq:s04:S} holds for $\XJ$ ($MJ\ge R+\frac12+\ell_c$); \cref{cor:s04:G1} gives $F^K_{k,J}\le F^K_{k,J+1}$ at every $k$ (the scheme multiplier is common, \cref{lem:s02:alg}(a)), and $k\to\infty$ gives $F^K_J\le F^K_{J+1}$ on $\Omega''_+$. Being nondecreasing and bounded there, $J\mapsto F^K_J$ converges on $\Omega''_+$, a non-empty open subset of the maximal totally real subspace $(\mathrm i\R)^n\times\R^{2p}$; $\cD''$ is convex, hence connected, and the Montel--Vitali fact (T3) stated before \cref{thm:s11:W} gives $F^K_J\to F^K_\infty$ locally uniformly on $\cD''$. At $w^+_j=w_j$, $w^-_j=-w_j$ the total scalar source is $\sum_jw_j(h_j^+-h_j^-)=\sum_jw_jh_j$, so $F^K_{k,J}(z,w,-w)=Z^{a,K}_{k,J}(\zeta)$ exactly at every $k$; hence $G^K_J\to G^K_\infty$ locally uniformly on $\cD'\defeq\{\zeta:|\Re w_j|<\delta_j\}$, which lies in $\Dtube$ because there $\norm{k_\pm}_\infty\le\norm k_\infty<\sum_j\delta_j\norm{h_j}_\infty<H/2$.

\emph{Step 9 (the whole tube).} The family $\{G^K_J\}_{J\ge J_1(R+1/2)}$ is locally uniformly bounded on $\Dtube$ by \eqref{eq:s13:WKlim}; by Montel's theorem every subsequence has a further subsequence converging locally uniformly to a holomorphic function, any two such limits agree on the open set $\cD'$ by Step~8, and $\Dtube$ is connected, so all subsequential limits coincide: $G^K_J\to G^K_\infty$ locally uniformly on $\Dtube$, $G^K_\infty$ is holomorphic, and \eqref{eq:s13:WKlim} passes to $G^K_\infty$. The finitely many indices $J_1(R)\le J<J_1(R+\frac12)$ do not affect the limit.

\emph{Step 10 (positivity and the cone).} For a real $\zeta\in\Dtube$ pick $\epsilon>0$ with $-\epsilon\zeta\in\Dtube$. The function $\varphi_J(t)\defeq\ln G^K_J(t\zeta)$ is finite and convex on $[-\epsilon,1]$ (a pointwise limit of the convex functions $\ln Z^{a,K}_{k,J}(t\zeta)$, \cref{lem:s02:alg}(d), with positive limits at real points) and $\varphi_J(0)=0$, so $\varphi_J(1)\ge-\varphi_J(-\epsilon)/\epsilon\ge-C_\epsilon/\epsilon$, $C_\epsilon$ the exponent of \eqref{eq:s13:WKlim} at $-\epsilon\zeta$; hence $G^K_\infty(\zeta)\ge e^{-C_\epsilon/\epsilon}>0$. At the cone points of $\Dtube$ ($z\in(\mathrm i\R)^n$, $w\in\R^p$, $k(\zeta)\ge0$), \cref{cor:s04:G1} at every $k$ and $k\to\infty$ give $G^K_J\le G^K_{J+1}\le G^K_\infty$.

\emph{Step 11 (the Schwinger bounds and \Ezz).} Put $S^K_{n,p}(f;h)\defeq\partial^{(n)}_z\partial^{(p)}_wG^K_\infty(0)$, $\sigma_0=H/4$, $K'=\kappa\pi^2+K_-$. Normalise $\norm{f_i}_2=1$ and let $\rho_j\defeq\sigma_0/(p\wnorm{h_j})$. On the polydisc $\{|w_j|\le\rho_j\}\times\C^n$ we have $\norm k_\infty\le\sum_j|\Re w_j|\norm{h_j}_\infty\le\sum_j\rho_j\wnorm{h_j}=\sigma_0<H/2\le\sigma_+/2$ (\cref{lem:s02:E}(E2)), so the polydisc lies in $\Dtube$, where $G^K_\infty$ exists with the bound \eqref{eq:s13:WKlim}, whose exponent is at most $\norm{g(z)}_2^2+\kappa\pi^2\wnorm{k_+}+K_-\wnorm{k_-}\le\norm{g(z)}_2^2+K'\wnorm k\le\norm{g(z)}_2^2+K'\sigma_0$ (E2). Cauchy's estimate for the mixed first derivative in the $w$'s at $w=0$, $z$ fixed, gives $|\partial^{(p)}_wG^K_\infty(z,0)|\le\prod_j\rho_j^{-1}\sup|G^K_\infty|\le(p/\sigma_0)^p\prod_j\wnorm{h_j}\,e^{K'\sigma_0}e^{\norm{g(z)}_2^2}$, and $p^p\le e^pp!$ (E4). The function $z\mapsto\partial^{(p)}_wG^K_\infty(z,0)$ is entire ($\C^n\times\{w\}\subset\Dtube$ for $|\Re w_j|$ small) and bounded by $A\,e^{\norm{g(z)}_2^2}$ with $A\defeq e^{K'\sigma_0}(e/\sigma_0)^pp!\prod_j\wnorm{h_j}$; on the polydisc of radius $r=(2n)^{-1/2}$ about $0$, $\norm{g(z)}_2\le\sum_i|\Re z_i|\le nr$, so Cauchy's estimate for the mixed first $z$-derivative gives $|S^K_{n,p}|\le A\,r^{-n}e^{n^2r^2}=A\,(2n)^{n/2}e^{n/2}\le A\,(\sqrt2e)^n(n!)^{1/2}$ by (E4). Restoring the normalisation by multilinearity,
\[
|S^K_{n,p}(f_1,\dots,f_n;h_1,\dots,h_p)|\le e^{K'\sigma_0}(\sqrt2e)^n(n!)^{1/2}(e/\sigma_0)^p\,p!\prod_i\norm{f_i}_2\prod_j\wnorm{h_j}.
\]
By (E2), $\norm{f_i}_2\le\frac\pi{2\sqrt2}\fournorm{f_i}$ and $\wnorm{h_j}\le\fournorm{h_j}$, and $(n!)^{1/2}p!\le(n+p)!$, so the right side is at most $\sigma_{n+p}\prod_i\fournorm{f_i}\prod_j\fournorm{h_j}$ with $\sigma_N=e^{K'\sigma_0}C_1^NN!$, $C_1=2\max(e\pi/2,e/\sigma_0)$; for complex test functions each argument is split into real and imaginary parts, which gives the factor $2^{n+p}\le2^N$ absorbed in $C_1$. The $S^K_{n,p}$ are therefore continuous multilinear forms in the norms $\fournorm\cdot$, extend to $\cS(\R^3;\Lambda^2_\C)^n\times\cS(\R^3;\C)^p$ and are tempered distributions, as in \cref{thm:s05:C}(c), and $C_1^N/(N!)^{\beta-1}\to0$ for every $\beta>1$: this is \Ezz\ with $s=4$ and every $\beta>1$, with the constants of \cref{thm:s13:main}C(i).
\end{proof}
